\section{What We Learned and the Revised Plan}\label{sec:plan}
The reproduction holds, but the paper mixes gene sets and lets the labels pick settings. Without labels the ranking can flip, mostly through the number of clusters each rule allows. For scale, correlation gives good candidates and principal components do not. The plan changes as follows.

\textbf{1. Stronger label-free rules as the main experiment.} A stability rule picks the setting whose clusters survive subsampling best, in the spirit of MultiK and chooseR \cite{liu2021multik,pattersoncross2021chooser}. In MICA's small-data mode it varies $k$. A significance rule merges clusters that a test cannot separate, as in sc-SHC \cite{grabski2023significance}. A transfer rule gives each method, on each dataset, the resolution with the best mean ARI on the other datasets, as Ullmann et al. recommend \cite{ullmann2023overoptimistic}. It needs MICA's graph mode, so it runs on Klein, Zeisel and PBMC14K. Outside the paper's benchmark, on the Allen subset, MICA's default gives 44 clusters for 34 classes, a first check on whether it was tuned. The silhouette rule stays, in one shared embedding and over the same range of cluster counts for both methods. The stability rule gives the final clustering, and each rule's settings are committed to a configuration file before scoring. All rules run on the ten benchmark sets, with all genes and with 2,000 variable genes, and on the Allen subset. We count a win for MICA if, under the stability rule, its median ARI is higher on at least 7 of the 11 sets and higher on average by 0.05 or more.

\textbf{2. Repeats for every result.} We do at least five runs of MICA and of Leiden per setting and report the spread, since MICA's label-free ARI on PBMC14K varied by 0.2 between runs.

\textbf{3. Scaling judged on subsets first.} Brute force on sampled cells stays the reference. A GPU version of the correlation step and the exact distance comes after November 4, and only if MICA wins under at least one label-free rule. We will judge it by recall and by agreement with MICA's own graph, without labels. The full 2.3-million-cell run is a stretch goal, attempted only if a 300,000-cell run predicts that it fits in a day and in 754 GB.

\textbf{4. Two optional checks,} if time allows: Leiden on the correlation neighbor graph, and a hub correction \cite{amblard2022hubness} on the mutual information graph.

\textbf{5. A narrower annotation part.} If the main experiments finish by November 4, we run one comparison on the PBMC14K and Zeisel clusters: marker genes alone against marker genes plus regulator activity from scMINER's own network and activity steps. Both arms query one fixed snapshot of CellMarker 2.0 and PanglaoDB and use the same local open-weight model at temperature 0. Two team members score each statement as supported or not by the retrieved text, and we report their agreement.
